\section{Discussion}
\label{sec:discussion}

{\system} demonstrates that composite visualization authoring can benefit from
treating explicit visual structure as the primary authoring reference. By
representing recognizable charts as reusable blocks and preserving composition
relationships as explicit structures, {\system} provides a middle ground
between predefined chart construction and low-level graphical specification.
Our evaluation suggests that this abstraction can improve compositional
understanding and preserve a consistent workflow across a broader range of
visualization structures.

\textbf{Lessons learned.}
Participants generally found recognizable chart blocks and explicit
composition relationships easy to understand, and the exploratory study showed
that block-based workflow can transfer across different chart
types, coordinate systems, and structural forms. These results support our
design choice of keeping visual forms and composition relationships explicit,
while allowing dimensional changes to be introduced when needed. Compared with
DI, this reduces the need to reason about fine-grained graphical and structural
levels; compared with Tableau, it places greater emphasis on making the
resulting visual form visible before further dimensional adaptation. The study
also suggests that variation in composable units can increase expressive
flexibility without requiring a different composition workflow.

\textbf{Limitations.} First, {\system} supports
composing visualizations from existing units but does not allow users to create
new units, which may constrain finer-grained designs and more open-ended
authoring. Second, our perceptual principles currently serve as design
guidelines rather than a systematically validated taxonomy. The expert-derived
families and implemented blocks therefore provide only an initial coverage of
the design space; additional composable structures and alternative block
boundaries may remain unexplored. 
Third, our current dimensional adaptation covers only a limited
set of transitions, primarily within chart families and through faceting.
More general dimensional changes may involve nested structures or transformations
across different families. Extending
these cases could support a broader recommendation space, where the system
suggests possible structural adaptations while leaving the final choice to the
author.



\section{Conclusion}
\label{sec:conclusion}

We presented {\system}, a block-based authoring environment for composite
visualizations that combines explicit chart units, layout-referenced
composition, and dimensional adaptation. A formative study informed the block
granularity and family organization, while case studies and user evaluations
showed that the resulting authoring logic supports diverse composite
visualizations through a consistent and learnable workflow. Overall, our work
suggests that treating visual structure as an explicit authoring reference can
bridge the gap between predefined chart construction and low-level graphical
specification.